\chapter{2026 年 8 月数据快照}
\markboth{数据快照}{数据快照}

\begin{ledetext}
把这些数字当刻度读，不要当结论读。三个月之后它们就不准了，但它们所处的结构会存在很久。
两年后回头翻这一页，比较哪些变了、哪些其实一直没变，是这份附录真正的用处。
\end{ledetext}

\section{利率：全球央行第一次不再同步}

\begin{tablewrap}
\begin{xltabular}{\linewidth}{>{\raggedright\arraybackslash}p{0.24\linewidth}Xr}
\toprule
\bthead{央行 / 品种} & \bthead{水平与动向} & \bthead{时点} \\
\midrule
\textbf{美国联邦基金利率} & 目标区间 3.50\%--3.75\%，2026 年已连续五次维持不变。主席凯文·沃什称通胀仍然高企，市场一度讨论加息可能 & 2026.07 \\
\textbf{欧洲央行} & 上调存款便利利率，为三年来首次加息 & 2026.06 \\
\textbf{日本银行} & 以 7 比 1 加息 25 个基点至 1.0\%，为 1995 年 9 月以来最高；国债购买规模继续压降 & 2026.06 \\
\textbf{中国 1 年期 LPR} & 3.0\%，连续多月持平 & 2026.08.20 \\
\textbf{中国 5 年期以上 LPR} & 3.5\%，连续多月持平 & 2026.08.20 \\
\textbf{中国 7 天期逆回购} & 1.4\%，主要政策利率 & 2026.08 \\
\textbf{中国商业银行净息差} & 1.41\%，历史低位，构成降息空间的现实约束 & 2026.Q2 \\
\textbf{美国 10 年期国债} & 收益率升至 2025 年 1 月以来最高，长端与短端明显分化 & 2026.08 \\
\bottomrule
\end{xltabular}
\end{tablewrap}

日本央行这一轮正常化的完整路径值得单独记一笔：2024 年 3 月结束负利率并终结收益率曲线控制，
政策利率由 $-0.1\%$ 上调至 $0\sim0.1\%$；2024 年 7 月加至 0.25\%；2025 年 1 月加至 0.5\%；
2025 年 12 月加至 0.75\%；2026 年 4 月因中东能源冲击暂停；2026 年 6 月重启加息至 1.0\%。
横跨二十七个月，负利率时代就此收尾。

\section{债务：353 万亿美元，和它的另一面}

\begin{statrow}{4}
  \statitem{353}{万亿美元}{全球债务总额，占全球 GDP 约 305\%（IIF，2026.03）}
  \statitem{40}{万亿美元}{美国国债规模首次突破，可流通部分逾 30 万亿（2026.08）}
  \statitem{2.3}{万亿美元}{全球私募信贷管理规模，2020 年以来翻倍（2025）}
  \statitem{3000}{亿美元}{全球稳定币市值量级（2026.08）}
\end{statrow}

\begin{databox}{必须同时记住的一句}
353 万亿美元的债务，同时也是 353 万亿美元的资产。讨论「债务太多」之前，先问一句：这笔账是谁欠谁的。
这就是第 5 章那副眼镜的用法。
\end{databox}

\section{储备：黄金第一次超过美国国债}

欧洲央行 2026 年 6 月发布的《欧元的国际角色》报告显示，截至 2025 年底，
全球央行储备资产的构成发生了一次结构性变化。

\begin{barfig}
  \barrow{黄金}{27}{27\%}
  \barrow{美国国债}{22}{22\%}
  \barrow{其他美元资产}{20}{20\%}
  \barrow{欧元}{15}{15\%}
\end{barfig}

黄金占比从 2024 年底的 20\% 升至 27\%，首次超过美国国债的 22\%；
美元计价资产合计仍占 42\%，仍是最大的单一币种。
全球央行黄金储备已超过 3.6 万吨，接近布雷顿森林体系鼎盛时期的约 3.8 万吨水平；
2025 年央行净购金 850 吨；金价在 2026 年 1 月一度接近每盎司 5600 美元。

一个容易被忽略的细节：稳定币发行商 Tether 在 2025 年成为单一最大黄金买家，购入超过 100 吨。

\section{市场与中国经济}

\begin{tablewrap}
\begin{xltabular}{\linewidth}{>{\raggedright\arraybackslash}p{0.26\linewidth}Xr}
\toprule
\bthead{项目} & \bthead{数值} & \bthead{时点} \\
\midrule
\textbf{上证指数} & 3986.30 点，当月上涨 4.02\%，两市成交约 2.13 万亿元 & 2026.08.31 \\
\textbf{深证成指} & 14015 点，当月上涨 3.21\% & 2026.08.31 \\
\textbf{创业板指} & 3438.68 点，当月上涨 2.83\% & 2026.08.31 \\
\textbf{日经 225} & 首次突破 70000 点 & 2026.06 \\
\textbf{中国上半年 GDP} & 同比增长 4.7\% & 2026.H1 \\
\textbf{中国公募基金规模} & 约 38 万亿元 & 2026 年初 \\
\textbf{被动投资} & 全球规模首次超越主动管理 & 2025--2026 \\
\textbf{数据中心资本开支} & 同比增长约 57\%，进入超级周期 & 2025 \\
\bottomrule
\end{xltabular}
\end{tablewrap}

\section{把这一页留到 2028 年}

\begin{takeaway}{使用说明}
两年后再翻开这一页，请按这个顺序看：先看哪些数字变了，再问一句为什么变；
然后看哪些结构没变。大概率的结果是，数字面目全非，而机制一条没换。
如果真的有某条机制变了，那才是值得停下来认真研究的事。
\end{takeaway}

\cleardoublepage